\section{PLANEJAMENTO, EXECUÇÃO E ANÁLISE DE TESTES}

Após as etapas de modelagem e desenvolvimento do aplicativo WeFIND, a fase de testes é fundamental para garantir o funcionamento correto, a estabilidade técnica e a facilidade de uso do aplicativo. Segundo Pressman e Maxim (2021), os testes de software servem para identificar falhas, validar os requisitos levantados e certificar que o produto atende às necessidades dos usuários.

Neste capítulo, são apresentados os testes funcionais de caixa-preta realizados sobre as regras de negócio, a avaliação de usabilidade feita com base no método \textit{System Usability Scale} (SUS) com 20 usuários convidados, e a análise dos tempos de resposta do sistema.

\subsection{Estratégia e Testes Funcionais de Caixa-Preta}

Os testes funcionais foram realizados utilizando a técnica de teste de caixa-preta, que avalia as entradas e saídas do sistema sem a necessidade de analisar o código-fonte interno. Para garantir a conformidade da aplicação com o planejamento do projeto, cada cenário de teste foi diretamente vinculado aos Requisitos Funcionais (RF) estabelecidos no Capítulo 4, testando tanto o comportamento esperado em situações normais quanto o tratamento de erros de preenchimento.

O Quadro 4 apresenta os casos de testes executados, os requisitos funcionais correspondentes e os resultados obtidos.

\vspace{0.3cm}
\begin{table}[!htb]
\centering
\phantomsection\label{quadro:testes}
\small
\setlength{\tabcolsep}{3.5pt}
\renewcommand{\arraystretch}{1.2}
\begin{tabular}{|c|>{\centering\arraybackslash}p{1.8cm}|>{\raggedright\arraybackslash}p{2.6cm}|>{\raggedright\arraybackslash}p{3.8cm}|>{\raggedright\arraybackslash}p{3.6cm}|c|}
\hline
\multicolumn{6}{|c|}{\textbf{Quadro 4: Matriz de Testes Funcionais do WeFIND}} \\ \hline
\textbf{ID} & \textbf{Requisito} & \textbf{Cenário de Teste} & \textbf{Procedimento Executado} & \textbf{Resultado Esperado} & \textbf{Situação} \\ \hline
CT01 & RF01, RF02 & Cadastro de usuário com dados válidos & Preenchimento de nome, e-mail válido e senha com mais de 6 caracteres & Criação da conta no Supabase Auth e direcionamento à tela inicial & Aprovado \\ \hline
CT02 & RF01 & Validação de e-mail duplicado & Tentativa de registro utilizando e-mail já cadastrado & Exibição de mensagem informativa e bloqueio do envio & Aprovado \\ \hline
CT03 & RF09, RF10 & Cadastro de animal com foto e mapa & Escolha de espécie, raça, ponto no mapa e foto da galeria & Gravação no banco, upload da imagem e exibição do marcador & Aprovado \\ \hline
CT04 & RF22, RF23 & Registro de avistamento no mapa & Marcação de ponto de avistamento com foto e relato rápido & Atualização do pino e alerta ao tutor responsável & Aprovado \\ \hline
CT05 & RF13, RF19 & Filtragem combinada na busca & Aplicação de filtro por espécie (Cão/Gato) e status (Perdido/Encontrado) & Atualização imediata da listagem e dos pontos no mapa & Aprovado \\ \hline
CT06 & RF37, RF38 & Envio de mensagens no chat & Troca de mensagens de texto entre duas contas conectadas & Entrega em tempo real e atualização direta na tela & Aprovado \\ \hline
CT07 & RF34, RF36 & Geração de cartaz com QR Code & Acionamento do comando ``Compartilhar Cartaz'' no detalhe do pet & Criação do cartaz com foto, dados e QR Code escaneável & Aprovado \\ \hline
CT08 & RF28, RF30 & Visualização do Perfil Digital do Pet & Preenchimento dos dados do pet e geração da tela de identificação & Exibição do Perfil Digital com dados do tutor e QR Code exclusivo & Aprovado \\ \hline
\end{tabular}
\fonte{Autoria própria (2026)}
\end{table}

Todos os casos de teste do Quadro 4 foram concluídos com sucesso, confirmando o pleno funcionamento das rotinas associadas aos requisitos funcionais analisados. As validações de campos responderam de forma clara nos casos de preenchimento incompleto, e as rotinas de cadastro, consulta e envio de arquivos mantiveram a integridade das informações no banco de dados.

\subsection{Avaliação de Usabilidade (Método SUS)}

Para avaliar a experiência de uso de forma prática e padronizada, foi aplicado o questionário \textit{System Usability Scale} (SUS), um método muito conhecido e validado na literatura para medir a facilidade de uso de sistemas digitais (LEWIS, 2021; SAURO; LEWIS, 2022). 

O questionário SUS conta com 10 afirmações respondidas em escala de 1 (discordo totalmente) a 5 (concordo totalmente), alternando perguntas positivas e negativas para evitar respostas automáticas (Apêndice B, Seção B.2). A pontuação final varia de 0 a 100 pontos, calculada pela fórmula:
\begin{equation}
    \text{Score SUS} = \left[ \sum (\text{itens ímpares} - 1) + \sum (5 - \text{itens pares}) \right] \times 2{,}5
\end{equation}

Os testes foram realizados com 20 voluntários (tutores de animais de estimação e pessoas interessadas na causa animal). Cada participante seguiu um roteiro com as seguintes tarefas:
\begin{enumerate}
    \item Criar uma conta de usuário e fazer login no aplicativo;
    \item Navegar pelo mapa interativo e abrir os detalhes de um animal cadastrado;
    \item Cadastrar um animal com foto e marcação de localização no mapa;
    \item Acessar a tela do pet cadastrado e gerar o cartaz com código QR;
    \item Acessar o Perfil Digital do Pet para visualização dos dados cadastrados.
\end{enumerate}

\subsubsection{Análise Quantitativa (Score SUS e Eficiência)}

Todos os 20 participantes conseguiram concluir as tarefas propostas sem necessidade de intervenção externa. O tempo médio total para percorrer todo o roteiro foi de \textbf{1 minuto e 45 segundos}. O preenchimento do formulário de cadastro com foto e marcação no mapa foi feito com agilidade, levando em média apenas \textbf{48 segundos}.

A média final das respostas dos 20 voluntários ao questionário SUS resultou em uma pontuação de \textbf{85,3 pontos}. De acordo com as faixas de classificação avaliadas por Lewis (2021) e Sauro e Lewis (2022), notas acima de 80 pontos colocam o sistema no nível de usabilidade \textbf{Excelente} (Grau A), confirmando que o WeFIND apresenta interface limpa, navegação intuitiva e é muito fácil de usar.

\subsubsection{Análise Qualitativa (Opinião dos Participantes)}

Após o teste prático, os voluntários compartilharam suas impressões sobre o aplicativo. Os comentários espontâneos mais frequentes destacaram:
\begin{itemize}
    \item \textbf{Visualização no mapa:} facilidade de enxergar os animais perdidos e encontrados diretamente na região onde o usuário mora ou costuma circular;
    \item \textbf{Cartaz automático com QR Code:} praticidade de salvar e imprimir um cartaz pronto sem precisar recorrer a editores de imagem;
    \item \textbf{Perfil Digital do Pet:} utilidade de centralizar foto, dados e contato de emergência em uma página exclusiva;
    \item \textbf{Privacidade e segurança:} segurança de poder trocar mensagens com outros usuários pelo chat interno sem expor o número pessoal de telefone.
\end{itemize}

\subsection{Desempenho e Tempos de Resposta}

O desempenho do WeFIND também foi avaliado em condições normais de conexão móvel (4G) e redes Wi-Fi residenciais:
\begin{itemize}
    \item \textbf{Carregamento do mapa e marcadores:} tempo médio de 1,2 segundo para posicionar todos os pontos visíveis na tela;
    \item \textbf{Envio de fotos:} tempo médio de 2,4 segundos para carregar uma nova imagem, auxiliado pela compressão automática que reduz o peso do arquivo antes do envio ao Supabase Storage;
    \item \textbf{Mensagens no chat:} entrega instantânea com tempo de propagação menor que 300 milissegundos via WebSockets.
\end{itemize}

Os resultados obtidos demonstram que o WeFIND atingiu estabilidade técnica, bom desempenho e facilidade de uso, cumprindo os objetivos propostos para o projeto.
\clearpage